%% kapa: exam=1 round=33 year=2022 period=2 form=A subject=accounting subject_ko=회계학 seq=39 number=79
%% source: Q-Net 감정평가사 기출문제 (cst00309) artlSeq=5209997 file=raw/1008/5209997/2022년도 제33회 감정평가사 1차시험 2교시(감정평가 관계법규, 회계학).pdf
%% conversion: kordoc 4.24.51 -> tools/gen1.py (auto); review=proofread-sonnet; flags=none
\begin{question}{79}
다음은 종합원가계산제도를 채택하고 있는 (주)감평의 당기 제조활동에 관한 자료이다.
\begin{framedtext}
기초재공품 \quad ￦3,000(300단위, 완성도 60\%) \par
당기투입원가 \quad ￦42,000 \par
당기완성품수량 \quad 800단위 \par
기말재공품 \quad 200단위(완성도 50\%) \par
\end{framedtext}
모든 원가는 공정 전체를 통하여 균등하게 발생하며, 기말재공품의 평가는 평균법을 사용하고 있다. 기말재공품원가는? (단, 공손 및 감손은 없다.)
\begin{choices}
  \item ￦4,200
  \item ￦4,500
  \item ￦5,000
  \item ￦8,400
  \item ￦9,000
\end{choices}
\end{question}
\answer{3}
